\section{Fixed-state Hugoniot fibres and their projections}
\label{sec:hugoniot}

Let $P\in\mathcal W$ and let
\[
  \pi_-(P)=U_-(P),\qquad \pi_+(P)=U_+(P).
\]
The symbols $U_-$ and $U_+$ denote local endpoint states attached to the moving point $P$; they are not global Riemann data.

\subsection{A generic fixed state}
For a generic exploration state $U_*\in\mathcal U$, neither a left nor a right role is intrinsic.  The two inverse images are
\[
 \mathcal H^{\mathrm{forw}}(U_*)=\pi_-^{-1}(U_*),\qquad
 \mathcal H^{\mathrm{back}}(U_*)=\pi_+^{-1}(U_*).
\]
Thus $U_-=U_*$ is fixed on the forward fibre while $U_+$ varies, and $U_+=U_*$ is fixed on the backward fibre while $U_-$ varies.

The full reciprocal transition used by these fibres is
\[
 Z=\frac1z,\qquad \eta=\frac{\tau}{Z^2},\qquad
 \Upsilon=\frac{Y}{Z}.
\]
Continuation through $Z=0$ is therefore a chart operation and must not be replaced by a renderer segment joining the two visual compactification boundaries.

\subsection{Characteristic lifts}
Intersections with $\mathcal C=\{Y=0\}$ are characteristic lifts of $U_*$.  When finite and spectrally distinct, the slow and fast lifts are denoted $P_s$ and $P_f$.  Coincidence and projective-infinity cases are retained as separate cases rather than forced into slow/fast labels.

\subsection{Spectral partition and strict Lax classification}
At every regular Hugoniot point, compute the ordered spectra
\[
 \lambda_1(U_-)\leq\lambda_2(U_-),\qquad
 \lambda_1(U_+)\leq\lambda_2(U_+)
\]
and the four gaps $s-\lambda_i(U_\pm)$.  For the strict Lax equivalence below,
both endpoint spectra must be strictly hyperbolic and all four gaps must be
nonzero.  Under these assumptions define
\[
 n_-=\#\{i:\lambda_i(U_-)<s\},\qquad
 n_+=\#\{i:\lambda_i(U_+)<s\}.
\]
For a $2\times2$ system, the strict Lax conditions used by the application are
\[
 (n_-,n_+)=(0,1)\quad\Longleftrightarrow\quad\text{Lax family 1},
\]
\[
 (n_-,n_+)=(1,2)\quad\Longleftrightarrow\quad\text{Lax family 2}.
\]
All other strictly hyperbolic signatures are labelled non-Lax only with respect to this strict Lax criterion.  No Liu condition or other admissibility criterion is implied.

Characteristic equalities and repeated endpoint eigenvalues require separate
analysis: counts alone do not enforce these exclusions.  The current
\texttt{laxClass2x2} implementation classifies counts without an independent
strictness check; its label on those boundaries is not a strict Lax certificate.

\subsection{Explicit projection to state space}
The two canonical projections must remain explicit.  Their action on the fixed-state fibres is
\[
\begin{array}{c|cc}
 &\pi_-&\pi_+\\ \hline
 \mathcal H^{\mathrm{forw}}(U_*)&\{U_*\}&\text{curve of }U_+\\
 \mathcal H^{\mathrm{back}}(U_*)&\text{curve of }U_-&\{U_*\}.
\end{array}
\]
Consequently the state-space view does not independently recompute a Hugoniot locus or its Lax class: it projects the already validated segments in $\mathcal W$.

\subsection{Riemann data and future saturation}
Global Riemann data have a different role:
\[
 U(x,0)=\begin{cases}U_L,&x<0,\\U_R,&x>0.\end{cases}
\]
The oriented inverse images are $\pi_-^{-1}(U_L)$ for waves leaving $U_L$ and $\pi_+^{-1}(U_R)$ for waves arriving at $U_R$.  This does not identify $U_L,U_R$ globally with the local symbols $U_-,U_+$.

For a moving curve $\Gamma\subset\mathcal W$, the local states may later seed the saturations
\[
 \operatorname{sat}_{\mathrm{forw}}(\Gamma)=
 \bigcup_{P\in\Gamma}\pi_-^{-1}(U_-(P)),\qquad
 \operatorname{sat}_{\mathrm{back}}(\Gamma)=
 \bigcup_{P\in\Gamma}\pi_+^{-1}(U_+(P)).
\]
These formulas define the next construction boundary; this Hugoniot closure does not yet construct composite waves or claim a general Riemann solver.
